\chapter{Introduction}\label{ch:intro}

\section{What quantum information is good for}
Quantum informatics, the use of quantum systems to store, transmit and process information, is
useful and important for
\begin{itemize}
  \item \textbf{basic science}: a better understanding of quantum physics, and the solution of quantum many-body problems;
  \item \textbf{communications}: quantum key distribution, quantum networks, quantum repeaters;
  \item \textbf{computing}: algorithms that are faster than any known classical algorithm for specific problems;
  \item \textbf{simulation}: using a controllable quantum system to emulate another one (this was the original motivation of Feynman in 1982);
  \item \textbf{sensing and metrology}: using entanglement and coherence to beat classical limits of precision.
\end{itemize}
Private companies are mostly interested in communications and computing, while academia can use
quantum computing and quantum simulation for research. Quantum machine learning, for example, can be useful for both.

A specialist in quantum computing in a company or in a research team should be able to answer the following questions:
\begin{itemize}
  \item Do we need a quantum computer at all? Is there a known quantum algorithm with an advantage for this problem?
  \item Can it work in our case? This requires an estimate of the error rate of the algorithm and of its implementation on a particular quantum computer
  (number of gates $\times$ error per gate, coherence time, connectivity).
  \item How do we map an optimization problem onto the hardware (QUBO, Ising model, embedding)?
\end{itemize}
In research, quantum computing is of interest for high-energy physics, nanophysics and photonics, i.e.\
the research fields of the laboratories of UCA (LPCA, Institut Pascal).

\section{Why quantum mechanics changes computation}
A classical bit is either $0$ or $1$. A \emph{qubit} is a quantum two-level system, whose state is a vector
\begin{equation}
  \ket\psi=\alpha\ket0+\beta\ket1,\qquad |\alpha|^2+|\beta|^2=1,
\end{equation}
in a two-dimensional complex Hilbert space. Three features of quantum mechanics, absent from classical
physics, are used by quantum computers:
\begin{enumerate}
  \item \textbf{Superposition and interference.} Amplitudes, not probabilities, are added. Amplitudes are complex
  numbers and can cancel each other (destructive interference) or reinforce each other (constructive interference).
  Quantum algorithms are clever arrangements of interference.
  \item \textbf{Entanglement.} The state of $n$ qubits lives in a space of dimension $2^n$, and a generic state cannot be
  described qubit by qubit. Describing it classically requires $2^n$ complex numbers: about $10^{15}$ bytes for $n=46$.
  This exponential size is both the reason why quantum systems are hard to simulate and the resource a quantum computer exploits.
  \item \textbf{Measurement.} The outcome of a measurement is random, and the measurement changes the state.
  At the end of a computation we obtain only $n$ classical bits out of the $2^n$ amplitudes. This severely limits
  what can be gained, and is the reason why only a few problems have a quantum speed-up.
\end{enumerate}

\section{Where we are}
Do not expect too much from this course: it is a general introduction to the topic. And do not expect too much from
quantum computing in the short term either. In spite of the ``Python interface'', we are at the age of ENIAC
(programming by hard-wiring registers), even before the age of the assembly language. At the same time the programming
paradigm is very different from classical programming, which makes things even harder. The positive point is that you can
become the new Knuth or Dijkstra.

The present era is often called NISQ (Noisy Intermediate-Scale Quantum, a term coined by J.~Preskill in 2018):
devices with $10^2$--$10^3$ physical qubits, error rates of $10^{-3}$--$10^{-2}$ per two-qubit gate, and no or only
rudimentary error correction. Error-corrected (``fault-tolerant'') quantum computers capable of running Shor's
algorithm on cryptographically relevant numbers would require of the order of $10^6$ physical qubits.

\section{Plan of the notes}
\begin{itemize}
  \item \Cref{ch:qm}--\cref{ch:applications1}: one qubit. Postulates of quantum mechanics, the two-level system and
  its physical origins, the Bloch sphere, measurement, and first applications (random numbers, no-cloning, BB84).
  \item \Cref{ch:entanglement}--\cref{ch:teleportation}: many qubits. Entanglement, Bell inequalities, the density
  matrix, decoherence and master equations, entanglement criteria, teleportation.
  \item \Cref{ch:gates}--\cref{ch:qec}: gates and algorithms. Universal gates, Qiskit, Deutsch--Jozsa, Grover, QFT and Shor,
  complexity, error correction.
  \item \Cref{ch:annealing}--\cref{ch:qml}: optimization, annealing, adiabatic computing and quantum machine learning.
  \item \Cref{ch:implementations}: physical implementations, with derivations of the Hamiltonians of superconducting
  qubits, trapped ions, cavity QED, NMR, spins in quantum dots, neutral atoms and photons.
\end{itemize}
